%=====================================================================
\chapter{Classic AI Systems and How to Analyse Them}\label{ch:classic}
%=====================================================================
\locator{6}{Part VI --- Learning, Language and Responsible AI}

\begin{outcomes}
By the end of this chapter you will be able to:
\begin{tight}
  \item \textbf{Describe} the General Problem Solver, Eliza, Student and Macsyma,
        and \textbf{identify} the single mechanism each one rests on.
  \item \textbf{Reconstruct} each of them in a few dozen lines using machinery from
        earlier chapters.
  \item \textbf{Apply} a six-column framework --- task, environment, representation,
        inference, knowledge source, failure mode --- to analyse any AI system.
  \item \textbf{Explain} what each system was mistaken for at the time, and why.
  \item \textbf{Use} the framework on a system you did not build.
\end{tight}
\end{outcomes}

\begin{prereq}
\chref{planning} (means--ends analysis), \chref{fol} (unification),
\chref{expert} (rule rewriting) and \chref{nlp} (pattern substitution). This chapter
introduces \emph{no new machinery}: everything it builds, you have already built.
\end{prereq}

\begin{keyterms}
General Problem Solver $\bullet$ Eliza $\bullet$ Student $\bullet$ Macsyma $\bullet$
means--ends analysis $\bullet$ pattern substitution $\bullet$ term rewriting
$\bullet$ microworld $\bullet$ the Eliza effect $\bullet$ system analysis $\bullet$
failure mode $\bullet$ knowledge source
\end{keyterms}

\section{Four Systems, Four Mechanisms}

The HEC 2023 outline names four programs as required case studies. They were built
between 1961 and 1968, they were each spectacular at the time, and each one is
essentially a single mechanism applied with taste.

That is the point of this chapter, and it is why it comes last. You can now build
all four --- not simplified versions, but the actual mechanisms --- because you
already have the parts.

\smallskip
\noindent\begin{longtable}{@{}L{2.4cm}L{1.5cm}L{4.4cm}L{6.2cm}@{}}
\toprule
\rowcolor{primary!15}
\textbf{System} & \textbf{Year} & \textbf{The mechanism} & \textbf{Where you built
it} \\
\midrule
\endhead
General Problem Solver & 1961 & Means--ends analysis: reduce the difference between
where you are and where you want to be & \chref{planning}, where it solved the
release-planning problem \\
Eliza & 1966 & Pattern substitution with pronoun reflection & \chref{nlp}, in about
thirty lines \\
Student & 1964 & Pattern matching from English into simultaneous equations, then
algebra & \chref{fol}'s unification, restricted to strings \\
Macsyma & 1968 & Term rewriting: apply algebraic identities as rules until nothing
applies & \chref{expert}'s forward chaining, over expressions \\
\bottomrule
\caption{The four required case studies, and the chapter of this handout that
already contains each mechanism.}\label{tab:classic}
\end{longtable}

\section{The General Problem Solver}

Newell and Simon's GPS (1961) was the first program explicitly intended to be
\emph{general}: a single solver, given a description of a domain, that would work on
any problem in it. Its engine was means--ends analysis --- find a difference between
the current state and the goal, select an operator that reduces that difference,
and recursively achieve that operator's preconditions.

\chref{planning} built exactly this, and measured it: on the four-component release
problem it found an optimal plan considering $187$ operators where progression
search expanded $1{,}163$ states.

\begin{conceptbox}
GPS was also a claim about \emph{people}. Newell and Simon collected think-aloud
protocols from human subjects solving puzzles and argued that the traces matched
what GPS did --- which places it in the \emph{thinking humanly} quadrant of
\dsref{ai-axes}, not the \emph{acting rationally} one this course otherwise
inhabits.

That matters for how it should be judged. As a psychological model it was a serious
and influential contribution. As a general problem solver it failed, and the reason
is the one \chref{planning} demonstrated: means--ends analysis is sensitive to goal
ordering and incomplete, and the generality was in the \emph{engine}, while all the
difficulty turned out to be in describing the domain.
\end{conceptbox}

\section{Eliza}

\chref{nlp} covered the mechanism and the effect. What belongs here is the analysis.

Eliza had no goal, no representation of the conversation, and no model of the person
it was talking to. It scored \emph{one} on \chref{introduction}'s five-point
checklist --- it has a kind of knowledge, in that the patterns are inspectable and a
non-programmer could edit them --- and zero on goal, search, adaptation and
explanation.

And it was taken for a therapist. That combination --- lowest possible sophistication,
highest possible impression --- is why it is the most instructive system in this
chapter, and why \chref{trends} returns to it.

\section{Student}

Bobrow's Student (1964) read algebra word problems written in English and solved
them. The example everyone quotes is about ages and uncles; the delivery equivalent
would be:

\begin{quote}\ttfamily\small
the sprint has twice as many bugs as stories .\\
the number of stories is 12 .\\
how many bugs does the sprint have ?
\end{quote}

Student worked in two stages. Patterns turned English fragments into equations ---
\texttt{twice as many X as Y} became $X = 2Y$, \texttt{the number of X is N} became
$X = N$ --- and then an ordinary equation solver did the algebra. Answer: $24$.

\begin{alertbox}[The microworld, and what it concealed]
Student worked beautifully inside a tiny fragment of English about arithmetic, and
that fragment is called a \term{microworld}. The assumption of the era --- stated
openly --- was that microworlds would be extended and joined until they covered
ordinary language.

They were not, and the reason is visible in the mechanism. Student's patterns are
brittle in a specific way: \texttt{twice as many bugs as stories} works, and
\texttt{the bug count is double the story count} does not, and neither does
\texttt{there are two bugs for every story}. Each paraphrase needs its own pattern,
the paraphrases are unbounded, and nothing in the approach makes the next one
cheaper than the last.

That is the same shape as \chref{expert}'s knowledge acquisition bottleneck, in a
different subject. Notice it recurring: it is the characteristic failure of
hand-built symbolic systems, and recognising it in a new project is worth more than
any individual technique in this book.
\end{alertbox}

\section{Macsyma}

Macsyma (from 1968, at MIT) did symbolic mathematics: it manipulated algebraic
expressions as structures rather than evaluating them numerically. It could
differentiate, integrate, factor and simplify, and it remained in serious scientific
use for decades. Its descendants are Maxima, Mathematica and SymPy.

The core is \term{term rewriting}: represent an expression as a tree, and repeatedly
apply identities as rewrite rules until none applies.

\begin{examplebox}[Rewrite rules for effort expressions]
\begin{tabular}{@{}L{6.0cm}L{8.6cm}@{}}
\toprule
\textbf{Rule} & \textbf{Reading} \\
\midrule
$x + 0 \to x$ & adding nothing changes nothing \\
$x \times 1 \to x$ & scaling by one changes nothing \\
$x \times 0 \to 0$ & scaling by zero gives nothing \\
$x + x \to 2x$ & collect like terms \\
$(a \times x) + (b \times x) \to (a+b) \times x$ & distributivity, backwards \\
\bottomrule
\end{tabular}

\smallskip
Applying these to $(3 \times e) + (2 \times e) + 0$ gives $5 \times e$. Each rule is
a fact about arithmetic that a person can check, the engine is
\chref{expert}'s forward chaining, and the conflict-resolution question is the same
one: when two rules apply, which fires?
\end{examplebox}

\section{A Framework for Analysing an AI System}

The HEC 2017 outline asks for \emph{analysis} of AI systems, not merely description.
Here is a framework that works on any of them, including ones you did not build and
whose internals you cannot see.

\dsfig{%
\begin{tikzpicture}[font=\scriptsize,
  col/.style={rounded corners=3pt, draw=#1, line width=1pt, fill=#1!8,
              text=#1!50!black, text width=2.05cm, align=flush left,
              inner sep=5pt, minimum height=3.3cm},
  hd/.style={rounded corners=2pt, fill=#1, text=white,
             font=\tiny\bfseries, inner sep=2.5pt, text width=2.02cm,
             align=center}]

\node[col=primary]   (c1) at (0,0)     {\\[2pt]What is it \emph{for}? State the
  performance measure, as in \chref{agents}.};
\node[col=secondary] (c2) at (2.45,0)  {\\[2pt]What does it act in? Observable?
  Static? Adversarial?};
\node[col=greenhl]   (c3) at (4.90,0)  {\\[2pt]How is the domain written down?
  This is usually the whole design.};
\node[col=purplehl]  (c4) at (7.35,0)  {\\[2pt]What does it do with the
  representation? Search, chain, rewrite, match.};
\node[col=tealhl]    (c5) at (9.80,0)  {\\[2pt]Where did the knowledge come
  from --- and who can correct it?};
\node[col=accent]    (c6) at (12.25,0) {\\[2pt]How does it fail? Gracefully or
  off a cliff? Does it know?};

\node[hd=primary]   at (c1.north) {TASK};
\node[hd=secondary] at (c2.north) {ENVIRON\-MENT};
\node[hd=greenhl]   at (c3.north) {REPRESEN\-TATION};
\node[hd=purplehl]  at (c4.north) {INFERENCE};
\node[hd=tealhl]    at (c5.north) {KNOWLEDGE\\SOURCE};
\node[hd=accent]    at (c6.north) {FAILURE\\MODE};

\node[rounded corners=3pt, fill=lightbg, draw=primary!30, line width=0.9pt,
      text=primary!70!black, font=\scriptsize, text width=13.8cm,
      align=flush left, inner sep=6pt] at (6.1,-2.35)
  {\textbf{The last two columns are the ones people skip, and they are the ones that
   predict what will happen.} \emph{Knowledge source} tells you whether the system
   can be maintained by the people who understand the domain, or only by its
   authors --- which decides whether it survives contact with a changing business.
   \emph{Failure mode} tells you what happens at the edge of competence, and in
   particular whether the system signals that it has reached one.};
\end{tikzpicture}}%
{Six columns for analysing an AI system. Applying them to a system you did not build
is the skill this chapter is for.}{analysis-framework}

\smallskip
{\small
% six narrow columns: let TeX stretch rather than overflow, and make
% hyphens affordable (the L type deliberately makes them expensive,
% which is right for two- and three-column tables and not for this)
\setlength{\emergencystretch}{3em}\hyphenpenalty=100\relax
% six columns means ten inter-column gaps; at the default 6pt they
% alone take 2.1cm of the 16.2cm text block
\setlength{\tabcolsep}{4pt}
\rowcolors{2}{white}{lightbg}
\noindent\begin{longtable}{@{}L{1.8cm}L{2.0cm}L{2.4cm}L{2.2cm}L{2.1cm}L{3.4cm}@{}}
\toprule
\rowcolor{primary!15}
\textbf{System} & \textbf{Task} & \textbf{Represen\-tation} &
\textbf{Inference} & \textbf{Knowledge} & \textbf{Failure mode} \\
\midrule
\endhead
\textbf{GPS} & Reach a goal & Operators: pre, add, delete & Means--ends &
Hand-written & Incomplete; goal order matters. Reports ``no plan'' on solvable
problems \\
\textbf{Eliza} & Converse & Patterns with variables & Match and substitute &
Hand-written & None detectable. It never fails --- it always has a reply, which is
the danger \\
\textbf{Student} & Solve a word problem & Patterns, then algebra & Match, then
solve & Hand-written & Brittle to paraphrase. Each new phrasing needs a new
pattern \\
\textbf{Macsyma} & Simplify & Expression trees & Rewrite to a fixed point &
Hand-written & Rule interactions; may loop, or stop at a form nobody
wanted \\
\bottomrule
\caption{The four systems through the framework. Note that all four share one
knowledge source --- a person, writing it down --- which is the single fact that
most determines their history.}\label{tab:classic-analysis}
\end{longtable}}

\begin{conceptbox}
Read down the \emph{knowledge source} column. Every one of the four was hand-built,
and the same limit appears in all four biographies: the system worked, extending it
required a person, and the extensions ran out before the domain did.

Now read the \emph{failure mode} column, and notice that Eliza's is the worst of
the four precisely because it has none. GPS reports failure. Student produces
nothing on an unmatched phrasing. Macsyma leaves an expression unsimplified. Eliza
always replies, so its user has no signal at all --- and that is the property that
\chref{trends} argues is shared by the systems now in production.
\end{conceptbox}

\begin{worked}[Analysing a system you did not build]
Apply the framework to something with no published internals: an IDE's code
completion.

\smallskip
\textbf{Task.} Suggest the next tokens a developer will type. The performance
measure is \emph{acceptance rate} --- how often a suggestion is taken --- which is
worth writing down, because it is a measure of the developer's behaviour, not of
whether the code was any good. Compare \chref{agents}: a measure of the agent's
activity rather than of the environment's state, and therefore gameable.

\smallskip
\textbf{Environment.} Partially observable --- it sees the file and perhaps the
project, not the developer's intent. Dynamic; the code changes as they type.
Single-agent. Episodic, approximately.

\smallskip
\textbf{Representation.} Unknown from the outside, and the framework lets us say
something anyway. If it completes an identifier defined three files away, it has a
symbol table --- an explicit representation. If it produces plausible code calling a
method that does not exist, it does not.

\smallskip
\textbf{Inference.} Test it. A completion that is always correct about types and
never invents a name is doing symbolic lookup. One that is usually right and
occasionally confident and wrong is doing something statistical. Most current tools
do both, and the two failure signatures are distinguishable from outside.

\smallskip
\textbf{Knowledge source.} The corpus it was trained on, plus the current project.
Neither is correctable by the developer, which is a real difference from every
system in Table~\ref{tab:classic-analysis}: nobody can fix a wrong suggestion, only
decline it.

\smallskip
\textbf{Failure mode.} This is the column that decides how to use it. If it fails by
suggesting nothing, it is safe and occasionally unhelpful. If it fails by suggesting
something fluent and wrong, it is Eliza with a compiler --- and the mitigation is
that the compiler and the tests are the failure detector the system does not have.

\smallskip
\textbf{What the analysis produced.} Not a verdict, but three testable questions and
one operational rule: never accept a completion that the type checker and the tests
have not seen. That rule came from the last column, which is why it is the one not
to skip.
\end{worked}

\begin{pitfall}
\begin{tight}
  \item \textbf{Describing a system instead of analysing it.} ``Eliza simulates a
        therapist'' is a description. The analysis is that it matches patterns, has
        no representation, and cannot fail detectably.
  \item \textbf{Skipping the knowledge-source column.} It predicts whether the
        system survives a change in the business, and all four classic systems
        share the same answer.
  \item \textbf{Skipping the failure-mode column.} A system that always produces an
        answer has the most dangerous failure mode there is.
  \item \textbf{Mistaking a microworld result for a general one.} Student worked on
        a fragment of English and the fragment did not extend. Ask what the
        demonstration excluded.
  \item \textbf{Judging GPS as a problem solver when it was offered as a
        psychological model.} Identify the quadrant of \dsref{ai-axes} first, then
        choose the evidence.
  \item \textbf{Assuming you cannot analyse a closed system.} Five of the six
        columns can be filled from the outside by testing, as the worked example
        does.
\end{tight}
\end{pitfall}

%---------------------------------------------------------------------
\begin{lab}[Four Systems in One File]
Each of the four reconstructed with machinery from earlier chapters, then the
analysis table printed from the code.
\begin{lstlisting}[language=Python]
"""Chapter 19 lab -- the four classic systems, reconstructed.

Nothing here is new. GPS is Chapter 9's means-ends analysis, Eliza is
Chapter 18's pattern substitution, Student is Chapter 12's unification
restricted to strings, and Macsyma is Chapter 13's forward chaining
over expression trees.
"""

import re

# ===================================================================
# 1. GPS (1961): means-ends analysis          -- as in Chapter 9
# ===================================================================
OPERATORS = [
    # (name, preconditions, add, delete)
    ("code", set(), {"coded"}, set()),
    ("review", {"coded"}, {"reviewed"}, set()),
    ("test", {"reviewed", "staging free"},
     {"tested", "staging busy"}, {"staging free"}),
    ("deploy", {"tested"}, {"deployed"}, set()),
    ("release staging", {"staging busy"},
     {"staging free"}, {"staging busy"}),
]


def gps(state, goals, depth=0, trace=None):
    """Achieve each goal by reducing differences. Chapter 9's engine."""
    if trace is None:
        trace = []
    if depth > 20:
        return None, trace
    for g in goals:
        if g in state:
            continue
        achieved = False
        for (name, pre, add, dele) in OPERATORS:
            if g not in add:
                continue
            got, trace = gps(state, sorted(pre), depth + 1, trace)
            if got is None:
                continue
            state = got
            if not pre <= state:
                continue
            state = (state - dele) | add
            trace.append(name)
            achieved = True
            break
        if not achieved:
            return None, trace
    return state, trace


# ===================================================================
# 2. Eliza (1966): pattern substitution       -- as in Chapter 18
# ===================================================================
REFLECT = {"my": "your", "i": "you", "me": "you", "am": "are"}
ELIZA_RULES = [
    (r"^i am (.*)$", "How long have you been {0}?"),
    (r"^my (\w+) keeps (.*)$", "Tell me more about your {0}."),
    (r"^because (.*)$", "Is that the real reason?"),
]


def eliza(line):
    low = line.lower().strip().rstrip(".!?")
    for pat, tmpl in ELIZA_RULES:
        m = re.match(pat, low)
        if m:
            groups = [" ".join(REFLECT.get(w, w) for w in g.split())
                      for g in m.groups()]
            return tmpl.format(*groups)
    return "Can you tell me more about that?"


# ===================================================================
# 3. Student (1964): English -> equations     -- Chapter 12's matching
# ===================================================================
STUDENT_PATTERNS = [
    # "the sprint has twice as many bugs as stories"
    (r"^the \w+ has twice as many (\w+) as (\w+)$",
     lambda m: (m.group(1), ("times", 2, m.group(2)))),
    # "the number of stories is 12"
    (r"^the number of (\w+) is (\d+)$",
     lambda m: (m.group(1), ("const", int(m.group(2))))),
]


def student(sentences, question):
    """Turn English into equations, then solve them."""
    eqs = {}
    for s in sentences:
        low = s.lower().strip().rstrip(".")
        for pat, build in STUDENT_PATTERNS:
            m = re.match(pat, low)
            if m:
                var, rhs = build(m)
                eqs[var] = rhs
                break
        else:
            return None, eqs, "no pattern matched: %r" % s

    def value(var, seen=None):
        seen = seen or set()
        if var in seen:
            return None                 # circular definition
        rhs = eqs.get(var)
        if rhs is None:
            return None
        if rhs[0] == "const":
            return rhs[1]
        if rhs[0] == "times":
            inner = value(rhs[2], seen | {var})
            return None if inner is None else rhs[1] * inner
        return None

    m = re.match(r"^how many (\w+) does the \w+ have$",
                 question.lower().strip().rstrip("?"))
    if not m:
        return None, eqs, "question not understood"
    return value(m.group(1)), eqs, None


# ===================================================================
# 4. Macsyma (1968): term rewriting          -- Chapter 13's chaining
# ===================================================================
# an expression is a nested tuple: ("+", a, b) or ("*", a, b),
# a number, or a string (a symbol)
def rewrite_once(e):
    """One pass of the identities. Returns (expr, rule name or None)."""
    if not isinstance(e, tuple):
        return e, None
    op, a, b = e
    if op == "+" and b == 0:
        return a, "x + 0 -> x"
    if op == "+" and a == 0:
        return b, "0 + x -> x"
    if op == "*" and b == 1:
        return a, "x * 1 -> x"
    if op == "*" and (a == 0 or b == 0):
        return 0, "x * 0 -> 0"
    if op == "+" and a == b:
        return ("*", 2, a), "x + x -> 2x"
    # (a*x) + (b*x) -> (a+b)*x
    if (op == "+" and isinstance(a, tuple) and isinstance(b, tuple)
            and a[0] == "*" and b[0] == "*" and a[2] == b[2]
            and isinstance(a[1], int) and isinstance(b[1], int)):
        return ("*", a[1] + b[1], a[2]), "(a*x)+(b*x) -> (a+b)*x"
    # recurse
    na, r = rewrite_once(a)
    if r:
        return (op, na, b), r
    nb, r = rewrite_once(b)
    if r:
        return (op, a, nb), r
    return e, None


def simplify(e, limit=25):
    """Forward chaining to a fixed point -- Chapter 13, over trees."""
    steps = []
    for _ in range(limit):
        e2, rule = rewrite_once(e)
        if rule is None:
            return e, steps
        steps.append(rule)
        e = e2
    return e, steps


def show(e):
    if not isinstance(e, tuple):
        return str(e)
    return "(%s %s %s)" % (show(e[1]), e[0], show(e[2]))


if __name__ == "__main__":
    print("=" * 62)
    print("1. GENERAL PROBLEM SOLVER (1961) -- means-ends analysis")
    print("=" * 62)
    final, plan = gps(frozenset({"staging free"}), ["deployed"])
    print("   goal: deployed")
    print("   plan:", " -> ".join(plan))
    assert plan == ["code", "review", "test", "deploy"]
    print("   This is Chapter 9's engine, unchanged.")

    print()
    print("=" * 62)
    print("2. ELIZA (1966) -- pattern substitution")
    print("=" * 62)
    for line in ("I am worried about the release",
                 "my deploy keeps failing",
                 "because the config is wrong",
                 "the quarterly numbers look fine"):
        print("   > %-34s %s" % (line, eliza(line)))
    print("   No representation of the conversation. It cannot fail,")
    print("   which is the most dangerous failure mode of the four.")

    print()
    print("=" * 62)
    print("3. STUDENT (1964) -- English into equations")
    print("=" * 62)
    facts = ["the sprint has twice as many bugs as stories",
             "the number of stories is 12"]
    q = "how many bugs does the sprint have?"
    ans, eqs, err = student(facts, q)
    for f in facts:
        print("   given: %s" % f)
    print("   equations:", eqs)
    print("   question: %s" % q)
    print("   answer:", ans)
    assert ans == 24

    print()
    print("   ... and the brittleness:")
    for para in ("the bug count is double the story count",
                 "there are two bugs for every story"):
        ans2, _, err2 = student([para, facts[1]], q)
        print("   %-42s -> %s" % (para, err2 or ans2))
        assert ans2 is None
    print("   Each paraphrase needs its own pattern, and the")
    print("   paraphrases are unbounded. This is Chapter 13's")
    print("   knowledge acquisition bottleneck, in another subject.")

    print()
    print("=" * 62)
    print("4. MACSYMA (1968) -- term rewriting")
    print("=" * 62)
    expr = ("+", ("+", ("*", 3, "e"), ("*", 2, "e")), 0)
    print("   before:", show(expr))
    out, steps = simplify(expr)
    print("   after: ", show(out))
    for s in steps:
        print("      applied  %s" % s)
    assert out == ("*", 5, "e")
    print("   Forward chaining to a fixed point, over trees instead")
    print("   of over facts. Chapter 13's engine, same conflict-")
    print("   resolution question: when two rules apply, which?")

    print()
    print("=" * 62)
    print("THE ANALYSIS")
    print("=" * 62)
    ANALYSIS = [
        ("GPS", "reach a goal", "operators", "means-ends",
         "hand-written", "incomplete; goal order"),
        ("Eliza", "converse", "patterns", "match+substitute",
         "hand-written", "NONE DETECTABLE"),
        ("Student", "solve a problem", "patterns+algebra", "match, solve",
         "hand-written", "brittle to paraphrase"),
        ("Macsyma", "simplify", "expression trees", "rewrite",
         "hand-written", "loops; wrong fixed point"),
    ]
    print("%-9s %-16s %-17s %-17s %-13s %s"
          % ("SYSTEM", "TASK", "REPRESENTATION", "INFERENCE",
             "KNOWLEDGE", "FAILURE MODE"))
    print("-" * 100)
    for row in ANALYSIS:
        print("%-9s %-16s %-17s %-17s %-13s %s" % row)

    sources = {r[4] for r in ANALYSIS}
    assert sources == {"hand-written"}
    print()
    print("Every one of the four has the same knowledge source, and")
    print("that single fact explains all four biographies: each")
    print("worked, each needed a person to extend it, and each ran")
    print("out of people before it ran out of domain.")
\end{lstlisting}
Then add a Student pattern for one of the two paraphrases that currently fail, and
notice how long it takes and how little it generalises. That experience --- not the
argument about it --- is what the chapter is for.
\end{lab}

%---------------------------------------------------------------------
\begin{checkpoint}
\begin{tightnum}
  \item Name the mechanism behind each of the four systems, and the chapter of this
        handout in which you built it.
  \item Eliza's failure mode is recorded as ``none detectable''. Explain why that is
        worse than GPS's ``reports failure on solvable problems'', and relate it to
        how a modern language system should be evaluated.
  \item Every one of the four has the same knowledge source. State it, and explain
        what it predicts about a hand-built symbolic system deployed in a changing
        business.
\end{tightnum}
\end{checkpoint}

%---------------------------------------------------------------------
\begin{chaptersummary}
\begin{tight}
  \item The four required case studies are each a single mechanism applied with
        taste: \term{GPS} is means--ends analysis (\chref{planning}), \term{Eliza}
        is pattern substitution (\chref{nlp}), \term{Student} is pattern matching
        into equations (\chref{fol}), \term{Macsyma} is term rewriting
        (\chref{expert}).
  \item GPS was offered as a model of human problem solving --- the \emph{thinking
        humanly} quadrant --- and should be judged as one. Its generality was in
        the engine; the difficulty was in describing the domain.
  \item Student worked inside a \term{microworld} and the microworld did not extend.
        Each paraphrase needed its own pattern, which is \chref{expert}'s knowledge
        acquisition bottleneck in another subject.
  \item Analyse a system on six columns: \emph{task, environment, representation,
        inference, knowledge source, failure mode}. The last two are the ones people
        skip and the ones that predict what will happen.
  \item All four share one knowledge source --- a person writing it down --- and
        that single fact explains all four biographies.
  \item Eliza's failure mode is the worst of the four because it has none: it always
        replies, so its user gets no signal. That property is shared by systems now
        in production, which is why \chref{trends} returns to it.
  \item Five of the six columns can be filled from outside a system you did not
        build, by testing. That is the transferable skill in this chapter.
\end{tight}
\end{chaptersummary}

\begin{reviewq}
  \item \textbf{(MCQ)} The General Problem Solver's engine was:
        (a)~resolution (b)~means--ends analysis (c)~term rewriting (d)~pattern
        substitution.
  \item \textbf{(MCQ)} Macsyma's core mechanism is: (a)~numerical evaluation
        (b)~term rewriting to a fixed point (c)~backward chaining (d)~unification
        with the occurs check.
  \item \textbf{(MCQ)} A \emph{microworld} is: (a)~a small computer
        (b)~a restricted domain in which a technique works, which was expected to
        extend (c)~a simulation environment (d)~a subset of first-order logic.
  \item \textbf{(MCQ)} Which failure mode is the most dangerous?
        (a)~reports failure on solvable problems (b)~brittle to paraphrase
        (c)~no detectable failure (d)~may loop.
  \item \textbf{(Short)} Explain why Student's brittleness and the expert-system
        knowledge acquisition bottleneck are the same phenomenon.
  \item \textbf{(Short)} GPS was a psychological model and a problem solver. Say how
        the evidence you would collect differs between those two claims.
  \item \textbf{(Short)} Give two of the six analysis columns that can be filled
        from outside a closed system, and describe the test you would run for each.
  \item \textbf{(Applied)} Apply the six-column framework to a spam filter, a
        satellite navigation system, or a tool you use daily. Fill every column,
        and state the one operational rule your analysis of the failure-mode column
        produces.
\end{reviewq}
